% META: {"id": "TEXP-ARI-CO-018", "chapitre": "TEXP-ARITHMETIQUE", "type_objet": "corrige", "exercice_ref": "TEXP-ARI-EX-018", "capacites_codes": ["C4"], "sources_inspiration": [], "status": "generated"}
% BEGIN-VERIFY
% import math
% assert math.gcd(7, 11) == 1
% assert 11 == 1 * 7 + 4 and 7 == 1 * 4 + 3 and 4 == 1 * 3 + 1
% assert 2 * 11 - 3 * 7 == 1
% assert 7 * (-3) + 11 * 2 == 1
% for k in range(-6, 7):
%     assert 7 * (-3 + 11 * k) + 11 * (2 - 7 * k) == 1
% trouvees = {(x, y) for x in range(-40, 41) for y in range(-40, 41) if 7 * x + 11 * y == 1}
% attendues = {(-3 + 11 * k, 2 - 7 * k) for k in range(-3, 4)}
% assert attendues <= trouvees
% END-VERIFY
\begin{corrige}{TEXP-ARI-EX-018}
\textbf{1.} $\mathrm{PGCD}(7\,;\,11)=1$ et $1$ divise $1$ : d'après le
théorème de Bézout, $(E)$ admet des solutions.

\textbf{2.} $11=1\times7+4$, $7=1\times4+3$, $4=1\times3+1$. En remontant :
$1=4-3=4-(7-4)=2\times4-7=2(11-7)-7=2\times11-3\times7$. Donc
$(x_0\,;\,y_0)=(-3\,;\,2)$ convient.

\textbf{3.} Si $(x\,;\,y)$ est solution, $7(x+3)=-11(y-2)$. Comme $7$ et $11$
sont premiers entre eux, $11$ divise $x+3$ : $x=-3+11k$, et alors $y=2-7k$.
Réciproquement ces couples conviennent. L'ensemble des solutions est
$\{(-3+11k\,;\,2-7k),\ k\in\mathbb{Z}\}$.
\end{corrige}
